% =============================================================================
% 2.4 本章小结
% 草稿版本：v0.1 (2026-05-05)
% 字数：约 500 字
% 风格规则：v0.2 风格（无 \textbf / \textit / 引例括号 / 双引号 / 破折号）
% 引用：暂无新增引用
% 依赖宏包：booktabs / tabularx（主稿已加载）
% =============================================================================

\subsection{本章小结}
\label{sec:problem-summary}

本章在第 1 章绪论给出的研究背景与文献综述基础上，把基于 AI 智能体的
气象预报数据检索与分析这一总体研究课题分解为三个相互独立又依次衔接
的研究问题，并给出了本文统一的方法论框架与领域适配层。

\ref{sec:problem-modeling} 节把课题中的技术难点按数据流向建模为
气象任务指令解析、气象数据异构特征、大语言模型数值与逻辑幻觉
三个具体问题。指令解析问题关注用户输入到结构化请求的转化，需要
处理参数结构化、表达不完整性、对抗鲁棒性三类子任务；异构数据
问题关注结构化请求到原始数据的取数与归一化，需要解决数据源、
数据粒度、数据要素三个层面的异构性；数值与逻辑幻觉问题关注原始
数据到自然语言决策报告的解读与表达，需要应对引用编造、分级编造、
数值漂移、范围错位、场景误判 5 类典型幻觉。三个问题在工程上构成
依赖关系，异构数据问题依赖指令解析输出的结构化意图，数值与逻辑
幻觉问题依赖工具调用返回的结果。

\ref{sec:react-paradigm} 节给出统一应对三类问题的方法论框架：
基于 ReAct 范式的单智能体思维链模型。本文在原始 ReAct 循环之上
增加了 NLU 前置、桥接子组件可调用、显式记忆三处工程改造，并通过
单智能体与多智能体的工程取舍对照论证了在响应延迟、状态同步、
调试可观测性三个维度选择单智能体的合理性。气象 Agent 的总体架构
划分为 NLU 前置层、ReAct 推理引擎层、工具与组件层、报告生成出口
四个层次，覆盖从用户输入到决策输出的完整流水线。

\ref{sec:rag-prompt} 节给出方法论框架在领域上的具体落地：基于
RAG 的领域知识注入与基于 Prompt 工程的工具调用约束。前者把
领域知识库划分为术语定义、分级标准、作业规范三类条目，强制要求
出处三元组加可信度标签加版本号的元数据完整性；后者把 117 行
\texttt{SYSTEM\_\bk{}PROMPT} 模块化为工具清单、工具使用规则、RAG 检索
规则、桥接调用约束、出行场景推理规则、输出规则 6 段，与
few-shot 示例形成抽象规则与具体范例的互补关系。

本章给出的问题建模、方法论框架与领域适配层是后续 3 至 5 章具体
方法落地的总体指南。第 3 章在 NLU 前置层与工具学习层落地指令解析
与异构数据接入两个研究问题，第 4 章在工具与组件层落地数值与逻辑
幻觉问题以及最终报告生成出口，第 5 章对四个层次进行模块层加集成
层的系统化实证。

